% Generated draft: requires mathematical and source review.
\documentclass[11pt]{article}
\usepackage[margin=25mm]{geometry}
\usepackage{fontspec}
\setmainfont{Latin Modern Roman}
\usepackage{amsmath,amssymb,mathtools,unicode-math}
\setmathfont{Latin Modern Math}
\usepackage{enumitem,needspace,xurl,hyperref}
\setlength{\parindent}{0pt}
\setlength{\parskip}{5pt}
\setlength{\emergencystretch}{3em}
\newcommand{\N}{\mathbb{N}}
\newcommand{\Z}{\mathbb{Z}}
\newcommand{\Q}{\mathbb{Q}}
\newcommand{\R}{\mathbb{R}}
\newcommand{\C}{\mathbb{C}}
\newcommand{\F}{\mathbb{F}}
\newcommand{\A}{\mathbb{A}}
\newcommand{\PP}{\mathbb{P}}
\newcommand{\E}{\mathbb{E}}
\newcommand{\eps}{\varepsilon}
\newcommand{\dd}{\,\mathrm{d}}
\begin{document}
\section*{UnsolvedMath ID: 20002819}
\textbf{Source problem code:} catalogue entry
\subsection{Definitions and mathematical statement}
\textbf{Problem 1.7.} Define superinduction and restriction beyond algebra groups. The source suggests automorphism-induced theories and characteristic subgroups, and Johnson--Smith's association-scheme induction work. It does not require a particular compatibility or character-preservation axiom. The aid distinguishes a known coherent-theory extension from a canonical projected construction for arbitrary theories.

For a finite group \(G\) with a supercharacter theory, let \(\operatorname{SCF}(G)\subseteq\operatorname{CF}(G)\) be its superclass-function space, with the normalized character inner product, and let \(P_G\) be orthogonal projection. Equivalently, \(P_G\) averages over each superclass. For arbitrary theories on \(H\leq G\), define
\[
 \operatorname{SRes}_H^G=P_H\operatorname{Res}_H^G,
 \qquad
 \operatorname{SInd}_H^G=P_G\operatorname{Ind}_H^G
\]
on the corresponding superclass spaces. The aid reports that these are adjoint. If \(K_g\) is the superclass of \(g\), the projected induction is exactly
\[
 \operatorname{SInd}_H^G f(g)=
 \frac{|G|}{|H|\,|K_g|}\sum_{h\in K_g\cap H}f(h).
\]
This construction always gives superclass functions, not automatically ordinary characters or a transitive functorial theory.

\textbf{Reported positivity and integrality distinction.} If \(X\) is a block of irreducible characters, use the canonical supercharacter \(\sigma_X=\sum_{\chi\in X}\chi(1)\chi\). For an ordinary character \(\psi=\sum_\chi m_\chi\chi\), its projection is
\[
 P_G\psi=\sum_Xc_X\sigma_X,\qquad
 c_X=\frac{\sum_{\chi\in X}m_\chi\chi(1)}
            {\sum_{\chi\in X}\chi(1)^2}\in\Q_{\geq0}.
\]
It remains an ordinary character precisely when \(c_X\chi(1)\) is an integer for every \(X\) and \(\chi\in X\). Nonnegative rational supercharacter coefficients alone do not establish integrality.

\textbf{Reported complete transitivity defect.} For theories on \(K\leq H\leq G\), the direct restriction minus its two-step projected version is, on \(\operatorname{SCF}(G)\),
\[
 D_{K,H,G}=P_K\operatorname{Res}_K^H(I-P_H)
                \operatorname{Res}_H^G.
\]
Thus projected restriction is transitive for this triple exactly when this restricted operator vanishes. By adjunction, the corresponding induction defect is its adjoint, so induction has the same criterion. The domain restriction matters: vanishing on the relevant superclass space, not on every arbitrary class function, is required. The aid reports an explicit nonzero example for \(C_8\supset C_4\supset C_2\), but does not supply the theories or test function in this summary; those data are not invented here.

\textbf{Reported automorphism-compatible case.} Let \(\Gamma\leq\operatorname{Aut}(G)\), take \(\operatorname{SCF}(G)=\operatorname{CF}(G)^\Gamma\), and give a \(\Gamma\)-invariant subgroup \(H\) the theory induced by the restricted action. Ordinary restriction and induction preserve these invariant class-function spaces, so both projections act trivially on their outputs and the projected maps reduce to the ordinary ones. Characteristic subgroups satisfy this invariance requirement. A chain with compatible restricted actions is therefore transitive; unrelated automorphism choices are not covered.

The aid attributes an earlier generalization of superinduction for coherent supercharacter theories on arbitrary finite groups to de Stavola. Exact coherence hypotheses and a full citation are absent. The projected construction answers a flexible linear-algebraic version, while character preservation, categorical realization and stronger compatibility requirements need separate hypotheses. All reported statements require independent verification.
\subsection{Short English statement}
Superclass averaging gives a canonical adjoint superrestriction/superinduction pair for arbitrary finite-group theories. It may lose character integrality and transitivity; the aid gives an exact leakage operator detecting the latter. Compatible automorphism-induced theories on invariant subgroups recover ordinary induction and restriction.
\subsection{Sources}
\begin{itemize}
\item Local UnsolvedMath registry, permanent ID \(20002819\), catalogue entry, \texttt{aim-representation-theory-notes.json}, \texttt{notes[90]}, \emph{Supercharacters and combinatorial Hopf algebras}, Problem 1.7: \url{https://aimath.org/WWN/supercharacters/supercharacters.pdf}.
\item Original remark: automorphism theories and characteristic subgroups; Johnson and Smith (1986), identified by the aid as the one-page note \emph{A Note on Character Induction in Association Schemes}. No journal details are supplied.
\item Supplied research summary names de Stavola's coherent-theory extension and gives projected adjoints, integrality and transitivity criteria, and a \(C_8,C_4,C_2\) example without its full data. Attempt 1: partial\_result, partially\_solved, rigor ``proved'', low novelty confidence. The machine-generated aid is not peer reviewed and requires independent verification before citation; no PDF visit is claimed.
\end{itemize}
\end{document}
